\chapter{Introduction}
\label{chap:introduction}

\section{Stable stratification as a computational problem}
\label{sec:intro-problem}

A stably stratified boundary layer is one in which density increases downward
fast enough that vertical displacement costs work. Nocturnal atmospheric layers
over land, the seasonal thermocline, and the shelf water beneath a fresh surface
plume are all in this state for most of their existence. Transport across such a
layer is what sets the surface energy balance in a weather model and the
nutrient supply in an ocean model, and it is the single largest source of
disagreement between operational parameterisations \citep{halvorsen2019closure,
brandt2021stableabl}.

The difficulty is not that stratified turbulence is hard to write down. The
Boussinesq equations are unambiguous. The difficulty is that stratification
introduces a second length scale that shrinks as the flow becomes more stable,
while the viscous scale shrinks only as the Reynolds number grows. The Ozmidov
scale
\begin{equation}
  L_O = \left( \frac{\varepsilon}{N^{3}} \right)^{1/2},
  \label{eq:ozmidov-intro}
\end{equation}
where $\varepsilon$ is the dissipation rate and $N$ the buoyancy frequency,
marks the largest eddy that can still overturn against buoyancy. Above $L_O$ the
motion is wave like and transports almost no scalar. Below it the motion is
turbulent and transports almost all of it. As the bulk Richardson number rises
towards its critical value, $L_O$ falls through the viscous scale and the active
region of the flow becomes a set of thin, intermittent, and spatially sparse
sheets embedded in a quiescent background.

A uniform grid sized for those sheets is a grid sized for the smallest feature
anywhere in the domain, applied everywhere. In the campaign reported in
Chapter~\ref{chap:results}, a uniform grid meeting the resolution requirement of
the most stable case would have required $3.1 \times 10^{10}$ sites, of which
fewer than four percent would have carried any scalar flux at all. That ratio,
and not any subtlety of the physics, is what makes the regime expensive.

\section{Why a kinetic scheme}
\label{sec:intro-kinetic}

Lattice Boltzmann methods solve a discrete kinetic equation whose hydrodynamic
limit is the Navier Stokes system. The collision step is local, the streaming
step is a shift, and the pressure field never requires a global solve. On
distributed memory hardware that locality is worth a great deal, because the
communication volume per time step scales with the surface of a subdomain rather
than with the iteration count of an elliptic solver \citep{feld2018latticeperf}.

The property matters more than usual here. An adaptive grid changes the
subdomain surface every time it is rebuilt, and an elliptic solver on a
frequently rebuilt hierarchy spends most of its time reconstructing coarse grid
operators. A kinetic scheme pays only for the interpolation of populations
across a refinement boundary, which is a local and cheaply parallel operation.

Against this, lattice methods have a reputation for accuracy problems at
refinement interfaces. The usual failure is that the non equilibrium part of the
population, which carries the stress, is not rescaled correctly when the
relaxation time changes with the grid spacing. The result is a spurious source of
entropy at every interface, small in a single interface test and fatal on a
hierarchy that is rebuilt every few hundred steps \citep{lombardi2020refine}.
Section~\ref{sec:method-coupling} shows that a rescaling exists which conserves
the second order moments of both populations exactly, and Section~\ref{sec:val-interface}
measures what happens without it.

\section{Contributions}
\label{sec:intro-contributions}

This dissertation makes the following contributions.

\begin{enumerate}
  \item A two population lattice kinetic formulation for buoyancy coupled scalar
    transport in which the momentum population uses a multiple relaxation time
    operator on the D3Q27 velocity set and the scalar population uses a single
    relaxation time operator on the reduced D3Q7 set. The construction is given
    in Chapter~\ref{chap:theory} and its Chapman Enskog expansion recovers the
    Boussinesq system to second order with a scalar diffusivity that is
    independent of the momentum relaxation rate.
  \item A refinement criterion keyed on the ratio of the local Ozmidov scale to
    the local grid spacing, developed in Section~\ref{sec:method-criterion}. The
    criterion refines where mixing is active rather than where shear is large,
    which in a stratified layer are different places.
  \item An interface coupling that conserves mass, momentum, and the second
    order non equilibrium moments across a refinement jump, with the proof in
    Section~\ref{sec:method-coupling} and the interface convergence study in
    Section~\ref{sec:val-interface}.
  \item A twelve case campaign across bulk Richardson numbers from 0 to 0.32,
    reported in Chapter~\ref{chap:results}, which establishes that the turbulent
    Prandtl number depends on the resolved fraction of the Ozmidov scale and not
    only on the gradient Richardson number, and a revised closure that carries
    this dependence.
\end{enumerate}

\section{Organisation}
\label{sec:intro-organisation}

Chapter~\ref{chap:theory} sets out the continuum problem, the kinetic
discretisation, and the moment structure that the rest of the work depends on.
Chapter~\ref{chap:method} describes the solver: the collision operators, the
refinement machinery, the interface coupling, and the parallel decomposition.
Chapter~\ref{chap:validation} establishes correctness against analytic solutions,
against a spectral reference for unstratified channel flow, and against the known
collapse of momentum flux at critical stratification. Chapter~\ref{chap:results}
presents the stratified campaign and the closure that follows from it.
Chapter~\ref{chap:conclusion} states what the work settles and what it does not.
Appendix~\ref{app:coefficients} tabulates the lattice coefficients and moment
transform used throughout, and Appendix~\ref{app:configuration} gives a complete
solver configuration for one production case so the runs can be reconstructed.
